\section{2024 Applied \& Computational Mathematics}

\begin{exercise}
\label{applied:2024:1}
Let $A \in \mathbb { R } ^ { n \times n }$ be a non-singular matrix. Let $u$ , $v \in \mathbb { R } ^ { n }$ be column vectors. Define the rank 1 perturbation $\stackrel { \triangledown } { \boldsymbol { A } } = \boldsymbol { A } + u \boldsymbol { v } ^ { T }$ .

(a) Derive a necessary and sufficient condition for $\widehat { A }$ to be invertible.   
(b) Let $x$ , $z$ and $b$ be column vectors in $\mathbb { R } ^ { n }$ . Suppose one can solve $A z = b$ with ${ \mathcal { O } } ( n )$ floating-point operations (flops). Under the conditions derived in(a), design an algorithm to solve ${ \widehat { A } } x = b$ with ${ \mathcal { O } } ( n )$ flops, and provide justification for your answer.
\end{exercise}

\begin{solution}
\textbf{(a) Condition for Invertibility:} 
By the Matrix Determinant Lemma, $\det(\widehat{A}) = \det(A + uv^T) = \det(A)(1 + v^T A^{-1} u)$.
Since $A$ is non-singular, $\det(A) \neq 0$. Thus, $\widehat{A}$ is invertible if and only if:
\[ 1 + v^T A^{-1} u \neq 0. \]

\textbf{(b) $\mathcal{O}(n)$ Algorithm:}
Using the Sherman-Morrison formula, if $\widehat{A}$ is invertible, then:
\[ \widehat{A}^{-1} = A^{-1} - \frac{A^{-1} u v^T A^{-1}}{1 + v^T A^{-1} u}. \]
To solve $\widehat{A}x = b$, we compute $x = \widehat{A}^{-1} b = A^{-1}b - \frac{A^{-1} u (v^T A^{-1} b)}{1 + v^T A^{-1} u}$.
\begin{enumerate}
    \item Solve $Az = b$ for $z$. ($O(n)$ flops)
    \item Solve $Aw = u$ for $w$. ($O(n)$ flops)
    \item Compute the scalar $\alpha = v^T z$. ($O(n)$ flops)
    \item Compute the scalar $\gamma = 1 + v^T w$. ($O(1)$ flops)
    \item Compute $x = z - \frac{\alpha}{\gamma} w$. ($O(n)$ flops)
\end{enumerate}
Total cost is $O(n)$, and the algorithm is justified by the Sherman-Morrison identity.
\end{solution}

\begin{exercise}
\label{applied:2024:2}
Consider the integral

\[
\int_ {0} ^ {\infty} f (x) d x
\]

where $f$ is continuous, $f ^ { \prime } ( 0 ) \neq 0$ , and $f ( x )$ decays like $x ^ { - 1 - \alpha }$ with $\alpha > 0$ in the limit $x \to \infty$ .

(a) Suppose you apply the equispaced composite trapezoid rule with $n$ subintervals to approximate

\[
\int_ {0} ^ {L} f (x) d x.
\]

What is the asymptotic error formula for the error in the limit $n  \infty$ with $L$ fixed?

(b) Suppose you consider the quadrature from (a) to be an approximation to the full integral from 0 to $\infty$ . How should $L$ increase with $n$ to optimize the asymptotic rate of total error decay? What is the rate of error decrease with this choice of L? 5

(c) Make the following change of variable $x = { \frac { L ( 1 + y ) } { 1 - y } }$ L(1 + y) , $y = { \frac { x - L } { x + L } }$ x − L in the original integral to obtain

\[
\int_ {- 1} ^ {1} F _ {L} (y) d y.
\]

Suppose you apply the equispaced composite trapezoid rule; what is the asymptotic error formula for fixed $L$ ?

(d) Depending on $\alpha$ , which method - domain truncation or change-of-variable - is preferable?
\end{exercise}

\begin{solution}
\textbf{(a)} For fixed $L$, the composite trapezoid rule error for $\int_0^L f(x) dx$ with $h = L/n$ is:
\[ E_{\text{trap}}(L, n) \approx -\frac{L h^2}{12} f''(\xi) = \mathcal{O}(h^2) = \mathcal{O}(L^2 n^{-2}). \]

\textbf{(b)} The full error is $E \approx | \int_L^\infty f(x) dx | + | E_{\text{trap}}(L, n) |$.
As $L \to \infty$, $\int_L^ \infty f(x) dx \sim \int_L^\infty x^{-1-\alpha} dx \sim L^{-\alpha}$.
Thus $E \sim L^{-\alpha} + L^2 n^{-2}$.
To balance the terms, set $L^{-\alpha} \sim L^2 n^{-2} \implies L^{\alpha+2} \sim n^2 \implies L \sim n^{\frac{2}{\alpha+2}}$.
The total error decay rate is $E \sim n^{-\frac{2\alpha}{\alpha+2}}$.

\textbf{(c)} The new integrand is $F_L(y) = f(L\frac{1+y}{1-y}) \frac{2L}{(1-y)^2}$.
As $y \to 1$, $x = L\frac{1+y}{1-y} \to \infty$. Given $f(x) \sim x^{-(1+\alpha)}$, we have:
\[ F_L(y) \sim \left( \frac{1+y}{1-y} \right)^{-(1+\alpha)} \frac{1}{(1-y)^2} \sim (1-y)^{1+\alpha-2} = (1-y)^{\alpha-1}. \]
Near $y=1$, the error of the trapezoid rule for an integrand with a singularity $(1-y)^{\beta}$ is known to be $\mathcal{O}(n^{-(1+\beta)})$.
For $\beta = \alpha-1$, the asymptotic error is $\mathcal{O}(n^{-\alpha})$.

\textbf{(d)} Comparison:
Domain truncation gives $\mathcal{O}(n^{-p})$ with $p = \frac{2\alpha}{\alpha+2} = 2 - \frac{4}{\alpha+2}$.
Change-of-variable gives $\mathcal{O}(n^{-\alpha})$.
As $\alpha \to \infty$, $p \to 2$ while $\alpha$ keeps increasing. Thus for large $\alpha$, change-of-variable is far superior.
Specifically, $\alpha > \frac{2\alpha}{\alpha+2}$ is always true for $\alpha > 0$ unless there are other smoothness constraints. However, for small $\alpha$, the singularity $(1-y)^{\alpha-1}$ becomes stronger, potentially limiting the method if $\alpha \leq 1$.
\end{solution}

\begin{exercise}
\label{applied:2024:3}
Consider the Chebyshev polynomial of the first kind

\[
T _ {n} (x) = \cos (n \theta), \quad x = \cos (\theta), \quad x \in [ - 1, 1 ].
\]

The Chebyshev polynomials of the second kind are defined as

\[
U _ {n} (x) = \frac {1}{n + 1} T ^ {\prime} (x), \quad n \geq 0.
\]

(a) Derive a recursive formula for computing $U _ { n } ( x )$ for all $n \geq 0$ .   
(b) Show that the Chebyshev polynomials of the second kind are orthogonal with respect to the inner product

\[
\langle f, g \rangle = \int_ {- 1} ^ {1} f (x) g (x) \sqrt {1 - x ^ {2}} d x
\]

(c) Derive the 2-point Gaussian Quadrature rule for the integral

\[
\int_ {- 1} ^ {1} f (x) \sqrt {1 - x ^ {2}} d x = \sum_ {j = 1} ^ {2} w _ {j} f (x _ {j})
\]
\end{exercise}

\begin{solution}
\textbf{(a) Recurrence:} 
Since $T_{n+1}(\cos\theta) = \cos((n+1)\theta)$, differentiating with respect to $x$ gives:
\[ T_{n+1}'(\cos\theta) (-\sin\theta) = -(n+1) \sin((n+1)\theta) \implies U_n(x) = \frac{\sin((n+1)\theta)}{\sin\theta}. \]
Using the identity $\sin((n+2)\theta) + \sin(n\theta) = 2 \sin((n+1)\theta) \cos\theta$, we have:
\[ (\sin\theta) U_{n+1}(x) + (\sin\theta) U_{n-1}(x) = 2 (\sin\theta) U_n(x) x \implies U_{n+1}(x) = 2x U_n(x) - U_{n-1}(x). \]
Initial conditions: $U_0 = 1$, $U_1 = 2x$.

\textbf{(b) Orthogonality:}
Let $x = \cos\theta, dx = -\sin\theta d\theta, \sqrt{1-x^2} = \sin\theta$.
\[ \int_{-1}^1 U_n(x) U_m(x) \sqrt{1-x^2} dx = \int_{\pi}^0 \frac{\sin((n+1)\theta)}{\sin\theta} \frac{\sin((m+1)\theta)}{\sin\theta} \sin\theta (-\sin\theta d\theta) \]
\[ = \int_0^\pi \sin((n+1)\theta) \sin((m+1)\theta) d\theta = \frac{\pi}{2} \delta_{nm}. \]

\textbf{(c) 2-Point Gaussian Quadrature:}
The nodes are the roots of $U_2(x) = 2x U_1 - U_0 = 4x^2 - 1 = 0 \implies x_1 = -1/2, x_2 = 1/2$.
The weights for a symmetric interval are $w_1 = w_2 = \frac{1}{2} \int_{-1}^1 \sqrt{1-x^2} dx = \frac{1}{2} \cdot \frac{\pi}{2} = \pi/4$.
The rule is: $\int_{-1}^1 f(x) \sqrt{1-x^2} dx \approx \frac{\pi}{4} [ f(-1/2) + f(1/2) ]$.
\end{solution}

\begin{exercise}
\label{applied:2024:4}
Consider the boundary value problem

\[
- \frac {d}{d x} \left(a (x) \frac {d u}{d x}\right) = f (x), \quad u (0) = u (1) = 0
\]

where $a ( x ) > \delta \geq 0$ is a bounded differentiable function in $[ 0 , 1 ]$ . We assume that, although $a ( x )$ is available, an expression for its derivative, $\frac { d a } { d x }$ is not available.

(a) Using finite differences and an equally spaced grid in $\lfloor 0 , 1 \rfloor$ , $x _ { l } = h l , l = 0 , \ldots , n$ and $h = 1 / n$ , we discretize the ODE to obtain a linear system of equations, yielding an $O ( h ^ { 2 } )$ approximation of the ODE. After the application of the boundary conditions, the resulting coefficient matrix of the linear system is an $( n - 1 ) \times ( n - 1 )$ tridiagonal matrix.

Provide a derivation and write down the resulting linear system (by giving the expressions of the elements).

(b) Utilizing all the information provided, find a disc in $\mathbb { C }$ , the smaller the better, that is guaranteed to contain all the eigenvalues of the linear system constructed in part (a).
\end{exercise}

\begin{solution}
\textbf{(a) Discretization:}
We use the centered difference scheme at $x_l$:
\[ - \frac{1}{h} \left[ \left( a \frac{du}{dx} \right)_{l+1/2} - \left( a \frac{du}{dx} \right)_{l-1/2} \right] \approx f_l \]
\[ - \frac{1}{h} \left[ a_{l+1/2} \frac{u_{l+1} - u_l}{h} - a_{l-1/2} \frac{u_l - u_{l-1}}{h} \right] = f_l \]
\[ \frac{1}{h^2} [ -a_{l+1/2} u_{l+1} + (a_{l+1/2} + a_{l-1/2}) u_l - a_{l-1/2} u_{l-1} ] = f_l. \]
For $l=1, \dots, n-1$, with $u_0 = u_n = 0$, we have $A u = f$ where:
\[ A_{l,l} = \frac{a_{l+1/2} + a_{l-1/2}}{h^2}, \quad A_{l,l-1} = -\frac{a_{l-1/2}}{h^2}, \quad A_{l,l+1} = -\frac{a_{l+1/2}}{h^2}. \]

\textbf{(b) Eigenvalue Bounds:}
By Gerschgorin's Circle Theorem, every eigenvalue $\lambda$ lies in the union of discs $D_l = \{ z \in \mathbb{C} \mid |z - A_{ll}| \leq R_l \}$ where $R_l = \sum_{j \neq l} |A_{lj}|$.
Here $R_l = \frac{a_{l+1/2} + a_{l-1/2}}{h^2} = A_{ll}$.
Thus each disc $D_l$ is centered at $A_{ll}$ and has radius $A_{ll}$, which means they all touch the origin and lie within the interval $[0, 2 A_{ll}]$.
Let $a_{\max} = \| a \|_\infty$. Then $A_{ll} \leq \frac{2 a_{\max}}{h^2}$.
The union of these discs is contained in a single disc centered at $\frac{2 a_{\max}}{h^2}$ with radius $\frac{2 a_{\max}}{h^2}$.
\[ \mathcal{D} = \left\{ z \in \mathbb{C} : \left| z - \frac{2 a_{\max}}{h^2} \right| \leq \frac{2 a_{\max}}{h^2} \right\}. \]
\end{solution}

\begin{exercise}
\label{applied:2024:5}
(a) Verify that the PDE

\[
u _ {t} = u _ {x x x}
\]

is well posed as an initial value problem.

(b) Consider solving it numerically using the scheme

\[
\frac {u (t + k , x) - u (t - k , x)}{2 k} = \frac {- \frac {1}{2} u (x - 2 h , t) + u (x - h , t) - u (x + h , t) + \frac {1}{2} u (x + 2 h , t)}{h}.
\]

Determine this scheme’s stability condition.
\end{exercise}

\begin{solution}
\textbf{(a) Well-posedness:}
Applying the Fourier transform in space: $\hat{u}_t = (ik)^3 \hat{u} = -ik^3 \hat{u}$.
The solution is $\hat{u}(k, t) = e^{-ik^3 t} \hat{u}(k, 0)$.
Since $|e^{-ik^3 t}| = 1$, we have $\|\hat{u}(\cdot, t)\|_{L^2} = \|\hat{u}(\cdot, 0)\|_{L^2}$. The $L^2$ norm is conserved, so the IVP is well-posed.

\textbf{(b) Stability:}
Assume the correct scale of the divisor is $h^3$ for consistency. Let $u_j^n = g^n e^{ij\theta}$.
Substituting into the scheme (Leapfrog step):
\[ \frac{g^{n+1} - g^{n-1}}{2k} = \frac{1}{h^3} \left[ -\frac{1}{2} e^{-2i\theta} + e^{-i\theta} - e^{i\theta} + \frac{1}{2} e^{2i\theta} \right] g^n \]
\[ g^2 - 1 = \frac{2k}{h^3} \left[ i \sin(2\theta) - 2 i \sin\theta \right] g. \]
Let $S = \frac{k}{h^3} ( \sin(2\theta) - 2 \sin\theta ) = \frac{2k}{h^3} \sin\theta (\cos\theta - 1)$.
The quadratic is $g^2 - 2 i S g - 1 = 0$, with roots $g = i S \pm \sqrt{1 - S^2}$.
For stability, we require $|g| = 1$, which implies $1 - S^2 \geq 0 \implies |S| \leq 1$.
The maximum of $|f(\theta)| = |2 \sin\theta (\cos\theta - 1)|$ occurs at $\cos\theta = -1/2$, where $|f(\theta)| = \frac{3\sqrt{3}}{2}$.
Thus stability requires $\frac{k}{h^3} \cdot \frac{3\sqrt{3}}{2} \leq 1 \implies k \leq \frac{2}{3\sqrt{3}} h^3 \approx 0.3849 h^3$.
\end{solution}

\begin{exercise}
\label{applied:2024:6}
Consider the diffusion equation

\[
\frac {\partial v}{\partial t} = \mu \frac {\partial^ {2} v}{\partial x ^ {2}}, \quad v (x, 0) = \phi (x), \quad \int_ {a} ^ {b} v (x, t) d x = 0
\]

with $x \in [ a , b ]$ and periodic boundary conditions. The solution is to be approximated using the central difference operator $L$ for the 1D Laplacian.

\[
L v _ {m} = \frac {v _ {m + 1} - 2 v _ {m} + v _ {m - 1}}{h ^ {2}},
\]

and the following two finite different approximations, (i) Forward-Euler

\[
v _ {n + 1} = v _ {n} + \mu k L v _ {n}, \tag {1}
\]

and (ii) Crank-Nicolson

\[
v _ {n + 1} = v _ {n} + \mu k \left(L v _ {n} + L v _ {n + 1}\right). \tag {2}
\]

Throughout, consider $[ a , b ] = [ 0 , 2 \pi ]$ and the finite difference stencil to have periodic boundary conditions on the spatial lattice $[ 0 , h , 2 h , \ldots , ( N - 1 ) h ]$ where $h = { \frac { 2 \pi } { N } }$ and $N$ is even.

(a) Determine the order of accuracy of the central difference operator $\mathit { L v }$ in approximating the second derivative $v _ { x x }$ .   
(b) Using $v _ { m } ^ { n } = \sum _ { l = 0 } ^ { N - 1 } \hat { v } _ { l } ^ { n } \exp \left( - i \frac { 2 \pi l m } { N } \right)$ vˆ give the updates $\hat { v } _ { l } ^ { n + 1 }$ in terms of $\hat { v } _ { l } ^ { n }$ for each of the methods, including the case $l = 0$ .   
(c) Give the solution for $v _ { m } ^ { \pi }$ for each method when the initial condition is $\phi ( m \Delta x ) =$ $( - 1 ) ^ { m }$ .   
(d) What are the stability constraints on the time step $k$ for each of the methods, if any, in equations (1) and (2)? Show there are either no constraints or express them in the form $k \leq F ( h , \mu )$ .
\end{exercise}

\begin{solution}
\textbf{(a) Accuracy:}
By Taylor expansion, $L v = v_{xx} + \frac{h^2}{12} v_{xxxx} + O(h^4)$, so the order of accuracy is $\mathcal{O}(h^2)$.

\textbf{(b) Growth Factors:}
Let $\hat{v}_l^{n+1} = g_l \hat{v}_l^n$. The eigenvalue of $L$ corresponding to the $l$-th mode is $\lambda_l = \frac{2(\cos \theta_l - 1)}{h^2} = -\frac{4}{h^2} \sin^2(\frac{\theta_l}{2})$, where $\theta_l = \frac{2\pi l}{N}$.
(i) FE: $g_{FE, l} = 1 + \mu k \lambda_l = 1 - \frac{4\mu k}{h^2} \sin^2(\frac{\theta_l}{2})$.
(ii) CN: $g_{CN, l} = \frac{1 + \mu k \lambda_l / 2}{1 - \mu k \lambda_l / 2} = \frac{1 - \frac{2\mu k}{h^2} \sin^2(\frac{\theta_l}{2})}{1 + \frac{2\mu k}{h^2} \sin^2(\frac{\theta_l}{2})}$.

\textbf{(c) Solution for $(-1)^m$:}
$(-1)^m = e^{im\pi}$, so $\theta = \pi$. $\sin^2(\pi/2) = 1$.
(i) $v_m^n = g_{FE}(\pi)^n (-1)^m = (1 - \frac{4\mu k}{h^2})^n (-1)^m$.
(ii) $v_m^n = g_{CN}(\pi)^n (-1)^m = \left( \frac{1 - 2\mu k / h^2}{1 + 2\mu k / h^2} \right)^n (-1)^m$.

\textbf{(d) Stability:}
For FE, $|g_{FE}| \leq 1 \implies -1 \leq 1 - \frac{4\mu k}{h^2} \sin^2(\theta/2) \leq 1$.
The left inequality gives $k \leq \frac{h^2}{2\mu}$ (for $\sin^2=1$).
For CN, $|g_{CN}| \leq 1$ is always true for any $k > 0$, so it is unconditionally stable.
\end{solution}
